\documentclass{article}

\usepackage{amsmath,amssymb}
\usepackage{xurl}
\usepackage[hidelinks]{hyperref}
\setlength{\emergencystretch}{2em}

% Shared commands for notes in docs/paper-gaps/.

\DeclareUnicodeCharacter{2080}{\ifmmode{}_{0}\else\textsubscript{0}\fi}
\DeclareUnicodeCharacter{2081}{\ifmmode{}_{1}\else\textsubscript{1}\fi}
\DeclareUnicodeCharacter{2082}{\ifmmode{}_{2}\else\textsubscript{2}\fi}
\DeclareUnicodeCharacter{2099}{\ifmmode{}_{n}\else\textsubscript{n}\fi}

\newcommand{\Lean}{\textsc{Lean}}
\ExplSyntaxOn
\NewDocumentCommand{\leanid}{m}{
  \ifmmode
    \textsf{\detokenize{#1}}
  \else
    \group_begin:
    \urlstyle{sf}
    \tl_set:Nn \l_tmpa_tl {#1}
    \tl_replace_all:Nnn \l_tmpa_tl {\_} {_}
    \exp_args:NV \path \l_tmpa_tl
    \group_end:
  \fi
}
\ExplSyntaxOff
\providecommand{\tr}{\operatorname{tr}}
\newcommand{\tnleanIssueBase}{https://github.com/LionSR/TNLean/issues/}
\newcommand{\tnleanPrBase}{https://github.com/LionSR/TNLean/pull/}
\newcommand{\tnleanDocBase}{https://sirui-lu.com/TNLean/docs/}
\newcommand{\ghissue}[1]{\href{\tnleanIssueBase#1}{GitHub issue \##1}}
\newcommand{\ghpr}[1]{\href{\tnleanPrBase#1}{PR \##1}}
\newcommand{\doclink}[1]{\href{\tnleanDocBase#1.html}{\path{#1}}}

% Dirac notation and the identity operator, as in blueprint/src/macros/common.tex.
% \providecommand keeps note-local definitions authoritative.
\usepackage{dsfont}
\providecommand{\ket}[1]{|#1\rangle}
\providecommand{\bra}[1]{\langle #1|}
\providecommand{\braket}[2]{\langle #1 | #2 \rangle}
\providecommand{\ketbra}[2]{|#1\rangle\!\langle #2|}
\providecommand{\Id}{\mathds{1}}
\providecommand{\one}{\mathds{1}}
\providecommand{\C}{\mathbb{C}}
\providecommand{\MN}[1]{M_{#1}(\C)}
\providecommand{\MD}{M_D(\C)}

% External referencing for paper sources.  The build helper writes sanitized
% auxiliary files under build/paper-aux/ and excludes duplicated source labels.
\usepackage{xr}

% Import a generated paper auxiliary file with a caller-supplied label prefix.
% The candidate paths support builds from docs/paper-gaps/, the repository root,
% and build/paper-aux/ itself.
\newcommand{\importpaperaux}[2]{%
  \IfFileExists{../../build/paper-aux/#2.aux}{%
    \externaldocument[#1]{../../build/paper-aux/#2}%
  }{%
    \IfFileExists{build/paper-aux/#2.aux}{%
      \externaldocument[#1]{build/paper-aux/#2}%
    }{%
      \IfFileExists{#2.aux}{%
        \externaldocument[#1]{#2}%
      }{}%
    }%
  }%
}

% General external reference macros
\newcommand{\paperref}[2]{\ref{#1-#2}}
\newcommand{\papereqref}[2]{(\ref{#1-#2})}

% CPSV16 source (1606.00608 / MPDO-22-12-17-2.tex) external document and shortcuts
\importpaperaux{cpsv16-}{1606.00608/MPDO-22-12-17-2-tnlean}
\newcommand{\cpsvref}[1]{\paperref{cpsv16}{#1}}
\newcommand{\cpsveqref}[1]{\papereqref{cpsv16}{#1}}

% Verdict marker: \gapnote{<kind>}{<status>}. Kind is the policy
% classification (clarification, local-correction, scope-restriction,
% unfaithful, false-source, open-gap); status is open, wip (elimination
% underway), resolved, or historical. Machine-read by the site index;
% prints a verdict line.
\newcommand{\gapnote}[2]{\par\noindent\textbf{Verdict:} #1 (#2).\par}


\title{Scope of the Formal Canonical-Form Reductions Relative to PGVWC07}
\author{}
\date{2026-05-12}

\begin{document}
\maketitle
\gapnote{scope-restriction}{resolved}

\section{The source theorem}

Let $A=(A_i)_i$ be an arbitrary translation-invariant matrix product state
(MPS) tensor on a finite ring, with bond dimension $D$. The
translation-invariant canonical-form theorem of P\'erez-Garc\'ia, Verstraete,
Wolf, and Cirac applies directly to $A$
\cite[Theorem~\texttt{Th:TIcanonical}, lines~742--763]{PerezGarcia2007MPSReps}.
It decomposes each matrix as
\begin{align}
    A_i
        &=
        \begin{pmatrix}
            \lambda_1 A_i^1 & 0 & 0 \\
            0 & \lambda_2 A_i^2 & 0 \\
            0 & 0 & \cdots
        \end{pmatrix},
    \label{eq:pgvwc07_ti_canonical_form_scope_block_decomposition}\\
    1
        &\geq \lambda_j>0.
    \label{eq:pgvwc07_ti_canonical_form_scope_positive_weights}
\end{align}
For every block $j$, the theorem gives a diagonal positive-definite matrix
$\Lambda^j$ such that
\begin{align}
    \sum_i A_i^j A_i^{j\dagger}
        &= \mathbf{1},
    \label{eq:pgvwc07_ti_canonical_form_scope_unital_block}\\
    \sum_i A_i^{j\dagger}\Lambda^j A_i^j
        &= \Lambda^j.
    \label{eq:pgvwc07_ti_canonical_form_scope_dual_fixed_point}
\end{align}
Define the transfer map of block $j$ by
\begin{align}
    \mathcal{E}_j(X)
        &:= \sum_i A_i^j X A_i^{j\dagger}.
    \label{eq:pgvwc07_ti_canonical_form_scope_transfer_map}
\end{align}
The identity is the only fixed point of $\mathcal{E}_j$ after the usual scalar
normalization of the fixed-point line. The resulting representation has bond
dimension at most $D$
\cite[Theorem~\texttt{Th:TIcanonical}, lines~742--763]{PerezGarcia2007MPSReps}.

The theorem does not assume injectivity, condition C1, blockwise primitivity, a
pre-existing block decomposition, strict separation of the weights, or one
block in each weight class. These properties appear elsewhere in the source as
later hypotheses, consequences, or genericity conditions. Injectivity is
established only after canonical form
\cite[Proposition~\texttt{prop-inj}, lines~911--950]{PerezGarcia2007MPSReps},
and generic uniqueness is a separate result
\cite[Theorem~\texttt{thm-uniq}, lines~1002--1015]{PerezGarcia2007MPSReps}.

\section{Faithful proof order}

The proof of Theorem~\texttt{Th:TIcanonical} in
Ref.~\cite{PerezGarcia2007MPSReps} follows a specific sequence. A faithful
formalization must derive the block structure in that order rather than assume
the hypotheses of later canonical-form results.

\begin{enumerate}
    \item Lines~765--770 normalize the spectral radius. If $X$ is a
          positive-definite fixed point, the proof uses the gauge
          \begin{align}
              B_i
                  &:= X^{-1/2}A_iX^{1/2},
              \label{eq:pgvwc07_ti_canonical_form_scope_full_rank_gauge}\\
              \sum_i B_iB_i^\dagger
                  &= \mathbf{1}.
              \label{eq:pgvwc07_ti_canonical_form_scope_full_rank_unitality}
          \end{align}

    \item Lines~771--783 treat a positive fixed point that is not full rank.
          Its support projection $P_R$ is derived from the spectral
          decomposition and the positivity contradiction in lines~775--783.
          The resulting invariant-subspace relation is
          \begin{align}
              A_iP_R
                  &= P_RA_iP_R.
              \label{eq:pgvwc07_ti_canonical_form_scope_support_invariance}
          \end{align}
          Thus $P_R$ is a conclusion, not an assumption from a prepared block
          decomposition.

    \item Lines~785--815 use
          (\ref{eq:pgvwc07_ti_canonical_form_scope_support_invariance}) to
          split the finite-ring trace into its restrictions to $R$ and
          $R^\perp$. Cyclicity moves the $P_R$ term through the product and
          eliminates the mixed contribution against $P_{R^\perp}$. This
          calculation gives the recursive block decomposition and shows that
          the bond dimension does not increase.

    \item Lines~816--826 repeat the split within each block. If a unital block
          has a non-scalar fixed point $X$, the proof shifts $X$ to obtain a
          singular positive fixed point:
          \begin{align}
              \mathbf{1}-\lambda_1^{-1}X
                  &\geq 0,
              \label{eq:pgvwc07_ti_canonical_form_scope_fixed_point_shift}\\
              \mathcal{E}(\mathbf{1}-\lambda_1^{-1}X)
                  &= \mathbf{1}-\lambda_1^{-1}X.
              \label{eq:pgvwc07_ti_canonical_form_scope_shifted_fixed_point}
          \end{align}
          This step forces each final fixed-point space to be the scalar line
          generated by $\mathbf{1}$.

    \item Lines~827--832 apply the analogous argument to the dual map. A
          unitary conjugation diagonalizes the resulting positive-definite
          fixed point and produces $\Lambda^j$ in
          (\ref{eq:pgvwc07_ti_canonical_form_scope_dual_fixed_point}).
\end{enumerate}

The source theorem is therefore neither a primitive-block theorem, an
injective-block theorem, nor a theorem about a prepared block family. The
invariant-projection lemma alone is also insufficient: the proof must derive
the support projection from a positive fixed point and then perform the
finite-ring trace split before invoking recursion.

\section{Normalization orientation}

The source uses the transfer map
(\ref{eq:pgvwc07_ti_canonical_form_scope_transfer_map}) and the unital
normalization
(\ref{eq:pgvwc07_ti_canonical_form_scope_unital_block}). The local
canonical-form chapter often uses the dual, trace-preserving orientation
\begin{align}
    \sum_i A_i^{j\dagger}A_i^j
        &= \mathbf{1}.
    \label{eq:pgvwc07_ti_canonical_form_scope_left_canonical_orientation}
\end{align}
The conventions correspond after passing to the adjoint transfer map or the
conjugate-transposed Kraus family. A theorem stated only in the left-canonical
orientation is not literally Theorem~\texttt{Th:TIcanonical} of
Ref.~\cite{PerezGarcia2007MPSReps} unless it also supplies this dualization or
an equivalent gauge conversion.

\section{Conditional formal reductions}

The checked primitive fixed-point reduction assumes injective blocks,
trace-preserving normalization in the chosen gauge, strict ordering by weight
modulus, nonzero blocks, positive block dimensions, and primitivity. The
peripheral-spectrum variant retains the injectivity, normalization,
weight-ordering, nonzero-weight, and positive-dimension assumptions and
requires every block transfer map to have peripheral spectrum $\{1\}$. The
normal-canonical-form reduction starts with a weighted block decomposition
already equivalent to the original state and assumes irreducibility,
trace-preserving normalization, primitivity, pairwise distinct weight moduli,
nonzero weights, and positive bond dimensions.

These are valid conditional reductions, but they do not formalize the
arbitrary-input theorem of Ref.~\cite{PerezGarcia2007MPSReps}. Their
hypotheses already contain substantial block structure. Citing them as the
source theorem would add assumptions to later uses of canonical-form
existence.

\section{Current proof-path audit}
\label{sec:pgvwc07_ti_canonical_form_scope_current_proof_path_audit}

The formal development contains the following components.

\begin{enumerate}
    \item
          \leanid{Kraus.spectralUnitalGauge_isUnital_of_map_eigenvector}
          formalizes the spectral-radius normalization in lines~765--769 of
          Theorem~\texttt{Th:TIcanonical} in
          Ref.~\cite{PerezGarcia2007MPSReps}. Given a positive-definite
          eigenvector $\rho$ with $\mathcal{E}_A(\rho)=r\rho$ and $r>0$, it
          constructs
          \begin{align}
              B_i
                  &:= r^{-1/2}\rho^{-1/2}A_i\rho^{1/2},
              \label{eq:pgvwc07_ti_canonical_form_scope_spectral_unital_gauge}\\
              \sum_i B_iB_i^\dagger
                  &= \mathbf{1}.
              \label{eq:pgvwc07_ti_canonical_form_scope_spectral_unital_result}
          \end{align}

    \item
          \leanid{Kraus.unitalGauge_isUnital_of_map_fixedPoint}
          formalizes the full-rank fixed-point gauge in lines~767--769. For a
          positive-definite fixed point $\rho$, it sets
          \begin{align}
              B_i
                  &:= \rho^{-1/2}A_i\rho^{1/2},
              \label{eq:pgvwc07_ti_canonical_form_scope_unital_fixed_point_gauge}\\
              \sum_i B_iB_i^\dagger
                  &= \mathbf{1}.
              \label{eq:pgvwc07_ti_canonical_form_scope_unital_fixed_point_result}
          \end{align}
          The companion theorem \leanid{MPSTensor.sameMPV_unitalGauge}
          proves that this similarity preserves finite-ring coefficients.

    \item
          \leanid{MPSTensor.exists_singular_posSemidef_fixedPoint_of_unital_nonScalar_fixedPoint}
          formalizes the fixed-point shift in lines~819--826. A Hermitian
          fixed point of a unital block that is not scalar yields the singular
          positive fixed point $\lambda_{\max}\mathbf{1}-X$. The companion
          theorem
          \leanid{MPSTensor.exists_twoBlock_decomp_of_unital_nonScalar_fixedPoint}
          combines the shift with support-projection splitting to produce a
          strict two-block decomposition with the same matrix-product-vector
          (MPV) family. The hypothesis is ``non-scalar'', not merely
          $X\neq\mathbf{1}$, because every scalar multiple of the identity is
          fixed by a unital linear map and cannot force a nontrivial split.

    \item
          \leanid{MPSTensor.fixed_eq_scalar_of_isIrreducibleTensor_unital}
          formalizes the endpoint of lines~816--826. Every fixed point of an
          irreducible unital block is a scalar multiple of $\mathbf{1}$. The
          proof uses the bridge between irreducible MPS tensors and irreducible
          maps together with
          \leanid{Kraus.fixed_eq_scalar_of_irreducible_unital}, which
          formalizes Theorem~6.6 of Ref.~\cite{Wolf2012Quantum}.

    \item
          \leanid{MPSTensor.exists_unitary_diag_posDef_adjointFixedPoint_of_unital_of_isIrreducibleTensor}
          formalizes the dual-map diagonalization in lines~827--832 for an
          irreducible unital block. It applies the trace-preserving
          diagonalization theorem to $i\mapsto A_i^\dagger$. Translating back
          gives $B_i=U^\dagger A_iU$ and a diagonal positive-definite
          $\Lambda$ satisfying
          \begin{align}
              \sum_i B_iB_i^\dagger
                  &= \mathbf{1},
              \label{eq:pgvwc07_ti_canonical_form_scope_dual_diag_unital}\\
              \sum_i B_i^\dagger\Lambda B_i
                  &= \Lambda.
              \label{eq:pgvwc07_ti_canonical_form_scope_dual_diag_fixed_point}
          \end{align}

    \item
          \leanid{MPSTensor.exists_unital_data_of_irreducible} formalizes the
          Perron--Frobenius unital-gauge step in lines~765--770 for one
          irreducible nonzero block. It applies the adjoint
          Perron--Frobenius theorem to the conjugate-transposed Kraus family,
          obtains a positive-definite eigenvector $\rho$, and constructs
          (\ref{eq:pgvwc07_ti_canonical_form_scope_spectral_unital_gauge}).
          This is a single-block input, not the recursive theorem.

    \item
          \leanid{MPSTensor.exists_pgvwc07_unital_dualDiag_data_of_irreducible}
          composes the single-block steps in lines~765--770 and 816--832. It
          constructs the Perron--Frobenius unital representative, proves that
          its fixed points are scalar, and diagonalizes a positive-definite
          dual fixed point after unitary conjugation. This closes the
          single-block composition gap, but it does not include recursive
          finite-ring decomposition, positive weights, or the total
          bond-dimension bound.

    \item
          \leanid{MPSTensor.exists_pgvwc07_unital_dualDiag_blockwise} applies
          the preceding result to a prepared family of nonzero irreducible
          summands with unit weights. It constructs positive spectral-radius
          weights, unital representatives, scalar fixed points, and diagonal
          positive-definite dual fixed points blockwise while preserving the
          weighted MPV family. It assumes that the irreducible decomposition
          has already been obtained.

    \item
          \leanid{MPSTensor.exists_pgvwc07_unital_dualDiag_from_arbitrary}
          composes a positive-length nonzero-block decomposition with the
          blockwise theorem. Starting from an arbitrary tensor, it removes
          all-zero blocks, which vanish on nonempty rings, and places the
          retained blocks in the unital orientation. It supplies positive
          weights, scalar fixed points, diagonal positive-definite dual fixed
          points, and
          \begin{align}
              \sum_k D_k
                  &\leq D.
              \label{eq:pgvwc07_ti_canonical_form_scope_bond_dimension_bound}
          \end{align}
          This bound follows from the direct-sum dimension identity, without
          an empty-word trace calculation.

    \item
          \leanid{MPSTensor.PGVWC07PositiveLengthWitness} and
          \leanid{MPSTensor.exists_pgvwc07_positiveLengthWitness} package the
          preceding theorem. Their fields record nonzero unital blocks,
          positive real weights, diagonal positive-definite dual fixed points,
          scalar transfer-map fixed points, positive block dimensions,
          positive-length MPV equality, and
          (\ref{eq:pgvwc07_ti_canonical_form_scope_bond_dimension_bound}).
          Positivity already implies nonvanishing of the weights. This witness
          is an equivalent formulation, not a new proof route.

    \item
          \leanid{MPSTensor.PGVWC07PositiveLengthWitness.exists_weight_normalization}
          formalizes the finite-family normalization in lines~765--766. For a
          nonempty witness, division by the largest positive weight produces
          weights $\nu_k$ such that
          \begin{align}
              \lvert\nu_k\rvert
                  &\leq 1,
              \label{eq:pgvwc07_ti_canonical_form_scope_normalized_weight_bound}\\
              \lvert\nu_{k_0}\rvert
                  &= 1
              \label{eq:pgvwc07_ti_canonical_form_scope_maximal_normalized_weight}
          \end{align}
          for some $k_0$. At length $N$, this normalization introduces a
          global factor $s^N$ in the MPV coefficients.

    \item
          \leanid{MPSTensor.PGVWC07PositiveLengthWitness.block_count_pos_of_exists_ne_zero_mpv}
          proves that the witness is nonempty whenever the original tensor has
          a nonzero positive-length MPV coefficient. An empty block family
          gives an empty sum, hence zero, at every positive length.

    \item
          \leanid{MPSTensor.exists_pgvwc07_normalized_exact_form_after_rescaling_of_exists_ne_zero_mpv}
          combines the arbitrary-input witness, its nonemptiness, and exact
          weight normalization. Under the nonzero positive-length MPV
          hypothesis, it produces $s>0$ and normalized weights satisfying
          (\ref{eq:pgvwc07_ti_canonical_form_scope_normalized_weight_bound})
          and
          (\ref{eq:pgvwc07_ti_canonical_form_scope_maximal_normalized_weight}).
          The tensor $s^{-1}A$ has the same positive-length MPV coefficients
          as the normalized weighted block tensor. This is the exact
          finite-ring effect of the ``without loss of generality''
          normalization in lines~765--766 of
          Ref.~\cite{PerezGarcia2007MPSReps}.

    \item
          \leanid{MPSTensor.exists_pgvwc07_normalized_exact_form_after_rescaling_or_forall_pos_mpv_eq_zero}
          gives the arbitrary-input dichotomy: either every positive-length
          MPV coefficient vanishes, or the exact rescaled normalized form
          exists.

    \item
          \leanid{MPSTensor.exists_pgvwc07_normalized_exact_form_after_rescaling_allow_empty}
          absorbs the zero branch into an empty retained block family. The
          positive-weight, unital, scalar-fixed-point, and dual-diagonal
          conditions are then vacuous. In the nonzero branch, the family is
          nonempty and contains a normalized weight of norm one. After a
          positive global rescaling, the block tensor agrees with the original
          tensor on every nonempty ring and satisfies
          (\ref{eq:pgvwc07_ti_canonical_form_scope_bond_dimension_bound}).

    \item
          \leanid{MPSTensor.exists_irreducible_blockDecomp} formalizes the
          recursive invariant-projection split in lines~765--833. It produces
          a finite direct sum of irreducible blocks with the same MPV
          coefficients and
          \begin{align}
              \sum_k D_k
                  &= D.
              \label{eq:pgvwc07_ti_canonical_form_scope_structural_dimension_identity}
          \end{align}
          This source-faithful partial result does not by itself construct the
          weights, unital orientation, diagonal dual fixed points, scalar
          fixed-point spaces, or final canonical-form package.

    \item
          \leanid{MPSTensor.exists_irreducible_blockDecomp_nonzeroBlocks}
          removes all-zero irreducible summands at positive lengths. It proves
          positive-length MPV equality, positive block dimensions, and
          (\ref{eq:pgvwc07_ti_canonical_form_scope_bond_dimension_bound})
          from
          (\ref{eq:pgvwc07_ti_canonical_form_scope_structural_dimension_identity}).

    \item
          \leanid{MPSTensor.exists_tp_gauge_from_arbitrary} applies
          Perron--Frobenius normalization blockwise after removing all-zero
          summands. It produces nonzero irreducible blocks satisfying
          \begin{align}
              \sum_i (B_k^i)^\dagger B_k^i
                  &= \mathbf{1}
              \label{eq:pgvwc07_ti_canonical_form_scope_arbitrary_tp_gauge}
          \end{align}
          and preserves the bond-dimension bound. This orientation is dual to
          the source's unital orientation and is equivalent only when the
          adjoint transfer map is carried explicitly.

    \item
          \leanid{MPSTensor.exists_normalCanonicalForm_of_primitive_blockDecomp}
          is a prepared-block theorem. It assumes a weighted block
          decomposition, irreducibility, trace-preserving normalization,
          primitivity, pairwise distinct weight moduli, nonzero weights, and
          positive bond dimensions. It remains a conditional corollary, not a
          citation target for Theorem~\texttt{Th:TIcanonical} of
          Ref.~\cite{PerezGarcia2007MPSReps}.
\end{enumerate}

The primitive, peripheral-primitive, and normal-canonical-form reductions are
not the source theorem. The blueprint instead uses the checked positive-length
statement as its formal counterpart of Theorem~\texttt{Th:TIcanonical}.
Dividing a finite positive family by its maximum gives the source convention
$1\geq\lambda_j>0$. Because the division changes every length-$N$ coefficient
by the factor $s^N$, the primary exact statement rescales the original tensor
by $s^{-1}$.

\section{Formal boundary}

The primitive, peripheral-primitive, and normal-canonical-form reductions must
be read as conditional canonical-form theorems, not as formalizations of the
arbitrary-input source theorem.\footnote{Section~\ref{sec:pgvwc07_ti_canonical_form_scope_current_proof_path_audit}
lists the formal declarations.}
They may be applied after source-faithful reductions establish their
hypotheses, but they may not be cited as producing canonical form directly
from an arbitrary translation-invariant MPS representation.

\section{Current conclusion}

The formal theorem used for Theorem~\texttt{Th:TIcanonical} of
Ref.~\cite{PerezGarcia2007MPSReps}, under the blueprint's positive-length
convention, is
\leanid{MPSTensor.exists_pgvwc07_normalized_exact_form_after_rescaling_allow_empty}.
It starts from an arbitrary translation-invariant MPS representation. After a
positive global rescaling, it constructs positive weights, unital blocks,
diagonal full-rank dual fixed points, scalar transfer-map fixed points, and the
bond-dimension bound
(\ref{eq:pgvwc07_ti_canonical_form_scope_bond_dimension_bound}). The retained
nonzero-block family is empty exactly when all positive-length MPV coefficients
vanish.

The preceding theorem
\leanid{MPSTensor.exists_pgvwc07_unital_dualDiag_from_arbitrary} directly
supplies the unnormalized positive-length block family and its bond-dimension
bound without a separate empty-word convention.

The primitive, peripheral-primitive, and normal-canonical-form reductions
remain later consequences available only after their additional hypotheses
have been proved. They are not alternative citations for
Theorem~\texttt{Th:TIcanonical}.

\bibliographystyle{alpha}
\bibliography{references}

\end{document}
